\section{Prinsip Induksi dan Analogi Domino}

Induksi matematika merupakan teknik pembuktian yang dirancang khusus untuk membuktikan pernyataan yang berlaku untuk seluruh bilangan bulat positif (atau bilangan bulat mulai dari suatu nilai tertentu). Berbeda dengan metode pembuktian pada Bab~12 yang bersifat umum, induksi matematika memanfaatkan struktur sekuensial bilangan bulat --- yaitu fakta bahwa setiap bilangan bulat memiliki penerus (\textit{successor}) --- untuk membangun bukti secara bertahap \cite{rosen2019}.

\subsection{Kapan Induksi Digunakan?}

Induksi matematika tepat digunakan ketika:
\begin{itemize}
    \item Pernyataan yang hendak dibuktikan berkaitan dengan bilangan bulat berurutan, misalnya ``untuk semua $n \geq 1$, $P(n)$ benar.''
    \item Domain melibatkan pola deret, formula penjumlahan, pertidaksamaan, atau sifat bilangan yang bergantung pada parameter bilangan bulat $n$.
    \item Verifikasi exhaustif tidak mungkin dilakukan karena domain tak berhingga.
\end{itemize}

\subsection{Analogi Efek Domino}

Prinsip induksi matematika dapat dipahami secara intuitif melalui analogi deretan batu domino yang berdiri berjajar. Untuk memastikan bahwa \emph{seluruh} batu akan jatuh, cukup dijamin dua hal:

\begin{enumerate}
    \item \textbf{Batu pertama dijatuhkan} (langkah basis): Kita secara aktif mendorong batu pertama sehingga ia jatuh.
    \item \textbf{Setiap batu yang jatuh akan meruntuhkan batu berikutnya} (langkah induksi): Jarak antar batu cukup dekat sehingga jika batu ke-$k$ jatuh, ia pasti menabrak dan menjatuhkan batu ke-$(k+1)$.
\end{enumerate}

Jika kedua kondisi ini terpenuhi, maka batu pertama jatuh (basis), menjatuhkan batu kedua (induksi $k=1$), yang menjatuhkan batu ketiga (induksi $k=2$), dan seterusnya tanpa batas.

\begin{figure}[H]
\centering
\begin{tikzpicture}[
    domino/.style={rectangle, draw=black!80, fill=blue!10, thick, minimum width=0.5cm, minimum height=1.5cm, rounded corners=1pt},
    fallen/.style={rectangle, draw=black!80, fill=red!20, thick, minimum width=0.5cm, minimum height=1.5cm, rounded corners=1pt, rotate=-60},
    arrow/.style={->, thick, >=stealth, red}
]
    % Fallen first domino
    \node[fallen] (d0) at (0,0) {};
    \node[below=0.6cm of d0, xshift=0.3cm] {$P(1)$};

    % Standing dominoes
    \node[domino] (d1) at (1.5,0) {};
    \node[below=0.2cm of d1] {$P(2)$};
    \node[domino] (d2) at (2.5,0) {};
    \node[below=0.2cm of d2] {$P(3)$};
    \node[domino] (d3) at (3.5,0) {};
    \node[below=0.2cm of d3] {$P(k)$};
    \node[domino] (d4) at (4.5,0) {};
    \node[below=0.2cm of d4] {$P(k\!+\!1)$};

    \draw[arrow] (0.6, 0.3) -- (1.2, 0.3);
    \draw[arrow] (d1.east) -- (d2.west);
    \node at (3,0) {$\ldots$};
    \draw[arrow] (d3.east) -- (d4.west);
    \node at (5.5,0) {$\ldots \; \infty$};

    % Labels
    \node[above=0.5cm of d0, font=\small\bfseries, text width=2cm, align=center] {Basis};
    \node[above=0.5cm of d3, font=\small\bfseries, text width=3.5cm, align=center, xshift=0.5cm] {Langkah Induksi};
\end{tikzpicture}
\caption{Ilustrasi Visual Efek Domino pada Induksi Matematika}
\label{fig:domino-induksi}
\end{figure}

\subsection{Pernyataan Formal Prinsip Induksi}

\begin{definisi}[Prinsip Induksi Matematika]
Misalkan $P(n)$ adalah pernyataan yang bergantung pada bilangan bulat $n$. Jika:
\begin{enumerate}
    \item $P(n_0)$ bernilai benar untuk suatu bilangan bulat $n_0$ (\textbf{langkah basis}), dan
    \item Untuk setiap bilangan bulat $k \geq n_0$, jika $P(k)$ benar maka $P(k+1)$ juga benar (\textbf{langkah induksi}),
\end{enumerate}
maka $P(n)$ bernilai benar untuk seluruh bilangan bulat $n \geq n_0$ \cite{wiki_mathematical_induction}.
\end{definisi}

Secara simbolis, prinsip induksi dapat ditulis sebagai:
\[ \Big[ P(n_0) \wedge \forall k \geq n_0\, (P(k) \rightarrow P(k+1)) \Big] \rightarrow \forall n \geq n_0\, P(n) \]

Perhatikan bahwa nilai awal $n_0$ tidak harus selalu $1$. Banyak teorema menggunakan $n_0 = 0$, $n_0 = 2$, atau bahkan nilai lainnya, tergantung pada domain pernyataan.
